% ========================================================================
% Lie Algebra Emergence as Neural Network Inductive Bias:
% A Systematic Experimental Campaign
%
% Author: Jose M. Moreno (independent researcher)
% Target: arXiv preprint
% ========================================================================
\documentclass[11pt,a4paper]{article}

% --- Packages ---
\usepackage[utf8]{inputenc}
\usepackage[T1]{fontenc}
\usepackage{amsmath,amssymb,amsthm}
\usepackage{mathtools}
\usepackage{physics}
\usepackage{booktabs}
\usepackage{multirow}
\usepackage{array}
\usepackage{graphicx}
\usepackage{xcolor}
\usepackage{hyperref}
\usepackage[margin=2.5cm]{geometry}
\usepackage{enumitem}
\usepackage{caption}
\usepackage{subcaption}
\usepackage{float}
\usepackage{tabularx}
\usepackage{authblk}

\hypersetup{
  colorlinks=true,
  linkcolor=blue!60!black,
  citecolor=green!50!black,
  urlcolor=blue!70!black
}

% --- Notation shortcuts ---
\newcommand{\fkg}{\mathfrak{g}}
\newcommand{\fksu}{\mathfrak{su}}
\newcommand{\fkso}{\mathfrak{so}}
\newcommand{\fksp}{\mathfrak{sp}}
\newcommand{\fksl}{\mathfrak{sl}}
\newcommand{\fku}{\mathfrak{u}}
\newcommand{\RR}{\mathbb{R}}
\newcommand{\CC}{\mathbb{C}}
\newcommand{\TK}{T_K}
\newcommand{\Tortho}{T_{\mathrm{ortho}}}
\newcommand{\Tclosure}{T_{\mathrm{closure}}}
\newcommand{\Tgate}{T_{\mathrm{gate}}}
\newcommand{\ngen}{n_{\mathrm{gen}}}
\newcommand{\hdual}{h^{\vee}}
\newcommand{\AlgNet}{\textsc{AlgNet}}
\newcommand{\GMLFixed}{\textsc{GML-Fixed}}

\theoremstyle{definition}
\newtheorem{definition}{Definition}
\newtheorem{observation}{Observation}

% ========================================================================
\title{%
  Lie Algebra Emergence as Neural Network Inductive Bias:\\
  A Systematic Experimental Campaign%
}

\author{Jose M. Moreno}
\affil{Independent Researcher}
\date{\today}

\begin{document}
\maketitle

% ========================================================================
% ABSTRACT
% ========================================================================
\begin{abstract}
We present a systematic experimental campaign (14 experiments, 18 generalizations) 
investigating whether Lie algebra structure can emerge as an inductive bias in neural 
networks through a variational pressure based on the Killing metric. The approach 
introduces an \emph{algebraic layer} whose parameters $\theta$ are constrained to the 
Lie algebra manifold $\fksu(d)$ via a basis expansion $G_a = \sum_k \theta_{ak} T_k$, 
and whose training includes a Killing tension term 
$\TK = \|K/\bar{K}_{\mathrm{diag}} + I\|_F^2$ that drives the learned generators 
toward a genuine Lie algebra basis. Across experiments spanning compact algebras 
$\fksu(d)$ for $d = 2$--$8$, sub-algebra recovery ($\fksu(3) \subset \fksu(4)$), 
non-compact algebras ($\fksl(3,\RR)$), and optimizer dynamics analysis 
(fractal dimension, Lyapunov exponents), we establish that: 
(i)~the algebraic parameterization is a \emph{necessary condition} for Killing 
crystallization; 
(ii)~a confinement term $\Tortho = \|\theta\theta^T - I\|_F^2$ is required to 
prevent manifold drift at $d \geq 4$; 
(iii)~a group-Lasso gate achieves perfect sub-algebra recovery; and 
(iv)~the decisive architectural variable is the \emph{decoder head}: a 2-layer 
head with width $h \geq 4\ngen$ unlocks the algebraic advantage 
($+6.65\%$ over MLP, $p \approx 0$, 10/10 seeds positive), while a 1-layer 
linear head creates an information bottleneck that masks the algebraic signal. 
The optimizer dynamics exhibit deterministic chaos 
($\lambda_\theta \approx +0.20$/step) with $\Tortho$ constraining the 
trajectory to a lower-dimensional fractal attractor ($D_2 = 0.99$ vs $1.41$ 
for MLP). These results establish concrete design rules for deploying algebraic 
inductive biases in neural networks and demonstrate that the Killing metric 
---~a purely geometric object~--- suffices as the sole variational driver of 
symmetry emergence in learned representations.
\end{abstract}

\tableofcontents

% ========================================================================
% 1. INTRODUCTION
% ========================================================================
\section{Introduction}
\label{sec:intro}

Symmetry is the organizing principle of modern physics. The Standard Model 
is built on the gauge group $SU(3) \times SU(2) \times U(1)$, general 
relativity on diffeomorphism invariance, and conformal field theory on the 
conformal group $SO(4,2)$. In each case, the symmetry group is 
\emph{postulated} as an axiom, and the dynamics are constructed to be 
compatible with it. A natural question arises: can symmetry instead 
\emph{emerge} from a simpler variational principle?

Recent work on the Killing metric variational principle~\cite{moreno2026relational,moreno2026experimental} 
has demonstrated that, given a set of $n = \dim(\fkg)$ random matrices 
subject to appropriate structural constraints (hermitian traceless for SU, 
antisymmetric for SO, symplectic for Sp), minimization of the Killing tension
\begin{equation}
  \label{eq:killing_tension}
  \TK = \left\| \frac{K}{\bar{K}_{\mathrm{diag}}} + I_n \right\|_F^2,
  \qquad K_{ab} = \sum_c \Tr([G_a, G_c][G_b, G_c]),
\end{equation}
drives convergence to an exact realization of the target Lie algebra 
$\fkg$ with 100\% probability across all tested algebras (SU($N$), SO($N$), 
Sp($2N$), G$_2$, F$_4$, $\fkso(3,1)$, $\fksu(2,2)$, and others). The dual 
Coxeter number $\hdual$ is recovered as an exact, parameter-free prediction 
for all 15 algebras tested.

In machine learning, the complementary question is: \emph{can algebraic 
structure serve as an inductive bias?} Equivariant neural networks 
\cite{cohen2016group,bronstein2021geometric,weiler2019general,fuchs2020se3,thomas2018tensor} 
hard-code known symmetry groups into the architecture, achieving 
state-of-the-art results on tasks with those exact symmetries. However, 
these approaches require the symmetry to be specified \emph{a priori}. 
Symmetry discovery methods~\cite{dehmamy2021automatic,yang2023generative,ling2022symmetry} 
aim to learn the symmetry from data, but typically operate on the group level 
rather than the algebra level.

This paper takes a different approach. We ask: \textbf{if we add the Killing 
tension~\eqref{eq:killing_tension} as a regularization term during neural 
network training, will the network spontaneously crystallize algebraic 
structure from the data, and does this improve learning?}

We report a systematic campaign of 14 experiments that trace the answer 
through three phases:
\begin{enumerate}[label=(\roman*)]
  \item \textbf{Manifold matters} (Exp.~1--3): Unconstrained matrices cannot 
    reach $\TK = 0$; the algebraic parameterization 
    $G_a = \sum_k \theta_{ak} T_k$ is a necessary condition.
  \item \textbf{Crystallization and extensions} (Exp.~4--9): On the correct 
    manifold, $\TK$ reliably crystallizes. Sub-algebra recovery via 
    group-Lasso, data efficiency, and non-compact algebra priors are 
    systematically tested.
  \item \textbf{Architecture as the hidden variable} (Exp.~10--14): The 
    algebraic layer \emph{always} crystallizes correctly (Killing signature 
    exact at all depths, dimensions, seeds). Whether crystallization 
    translates into a task advantage depends entirely on the decoder head: 
    a 2-layer head with width $h \geq 4\ngen$ is both necessary and 
    sufficient.
\end{enumerate}

The campaign yields 18 generalizations (G1--G18), organized thematically 
in Section~\ref{sec:generalizations}, and a set of concrete design rules 
for practitioners (Section~\ref{sec:design_rules}).


% ========================================================================
% 2. METHOD
% ========================================================================
\section{Method}
\label{sec:method}

% --- 2.1 ---
\subsection{Algebraic parameterization}
\label{sec:algebraic_param}

Let $\{T_k\}_{k=1}^{d^2-1}$ be a fixed orthonormal basis of $\fksu(d)$ 
(e.g., generalized Gell-Mann matrices), satisfying 
$\Tr(T_j T_k) = \tfrac{1}{2}\delta_{jk}$. We define the \emph{algebraic 
layer} as a set of $\ngen$ learned generators:
\begin{equation}
  \label{eq:algebraic_layer}
  G_a = \sum_{k=1}^{\ngen} \theta_{ak}\, T_k, 
  \qquad \theta \in \RR^{\ngen \times \ngen},
\end{equation}
where $\theta$ is a free parameter matrix ($\ngen = d^2 - 1$ for the 
full algebra). By construction, every $G_a$ is hermitian and traceless, 
hence $G_a \in \fksu(d)$ regardless of $\theta$. When 
$\theta \in O(\ngen)$ (orthogonal group), the set $\{G_a\}$ spans the 
same algebra as $\{T_k\}$.

The forward pass of the algebraic layer on an input 
$M(x) = \sum_b x_b T_b \in \fksu(d)$ is:
\begin{equation}
  \label{eq:forward}
  f_a(x) = \operatorname{Re}\Tr(G_a \cdot M(x)) 
  = 2 \sum_b \theta_{ab}\, x_b = 2(\theta\, x)_a.
\end{equation}
This is a bilinear projection that extracts algebraically structured 
features. When $\theta = I$ (identity), the layer is an exact coordinate 
extractor; when $\TK = 0$, $\theta \in O(\ngen)$ and the layer is an 
isometric rotation---no information is lost.

% --- 2.2 ---
\subsection{Loss functions}
\label{sec:losses}

The total loss combines a task-specific cross-entropy with up to three 
algebraic regularizers:
\begin{equation}
  \label{eq:total_loss}
  \mathcal{L} = \mathcal{L}_{\mathrm{CE}} 
  + \lambda_K\, \TK 
  + \lambda_{\mathrm{ortho}}\, \Tortho 
  + \lambda_{\mathrm{gate}}\, \Tgate.
\end{equation}

\paragraph{Killing tension $\TK$.}
Defined in~\eqref{eq:killing_tension}. Measures deviation of the Killing 
form $K_{ab}$ from the target $K \propto -I$ (for compact semisimple 
algebras). Reaches its global minimum $\TK = 0$ when and only when 
$\{G_a\}$ spans a compact semisimple Lie algebra.

\paragraph{Orthogonality tension $\Tortho$.}
\begin{equation}
  \Tortho = \|\theta\,\theta^T - I_{\ngen}\|_F^2.
\end{equation}
Confines $\theta$ to the Stiefel manifold $V_{\ngen}(\RR^{\ngen})$. 
Without this term, the task gradient pushes $\theta$ off the orthogonal 
manifold, with $\Tortho$ growing catastrophically at $d \geq 4$ 
(e.g., $\Tortho = 53$ for $d = 5$, Exp.~5E). This is the neural analog 
of the closure tension $\gamma \cdot \Tclosure$ used in lattice 
Monte~Carlo verification of exceptional algebras~\cite{moreno2026experimental}.

\paragraph{Gate tension $\Tgate$.}
For sub-algebra discovery, a per-generator scale 
$s_a = \operatorname{softplus}(\log s_a) > 0$ modulates each generator:
\begin{equation}
  G_a^{\mathrm{active}} = s_a \cdot G_a, 
  \qquad \Tgate = \sum_a s_a.
\end{equation}
This group-Lasso penalty~\cite{yuan2006model} drives inactive generators 
to zero while preserving the active sub-algebra.

\paragraph{Normalized Killing coefficient.}
The effective Killing gradient scales as $O(\lambda_K \cdot \ngen)$. 
To maintain constant per-generator pressure across dimensions:
\begin{equation}
  \label{eq:lambda_norm}
  \lambda_K^{\mathrm{eff}}(d) 
  = \lambda_K^{\mathrm{ref}} \cdot \frac{\ngen^{\mathrm{ref}}}{\ngen(d)}
  = 2.0 \cdot \frac{8}{d^2 - 1}.
\end{equation}

% --- 2.3 ---
\subsection{Task design}
\label{sec:task}

The benchmark task throughout the campaign is the \textbf{cubic invariant 
classification}:
\begin{equation}
  \label{eq:task}
  y(x) = \operatorname{sign}\!\left(\Tr(M(x)^3)\right), 
  \quad M(x) = \sum_a x_a T_a \in \fksu(d).
\end{equation}
This task is exactly invariant under the adjoint action of $SU(d)$: 
$\Tr(M^3) = \Tr((UMU^\dagger)^3)$ for all $U \in SU(d)$. It therefore 
has genuine algebraic symmetry, and any representation that respects the 
$SU(d)$ structure should achieve higher sample efficiency than an 
unconstrained model.

Key properties:
\begin{itemize}[nosep]
  \item Inputs: $x \in \RR^{d^2-1}$ (Gell-Mann coordinates of a random 
    $\fksu(d)$ element).
  \item Labels: binary ($\pm 1$), approximately balanced.
  \item Difficulty increases with $d$: accuracy ceiling drops toward 50\% 
    as $\ngen$ grows.
  \item For $d = 2$: $\Tr(M^3) = 0$ identically (all traceless $2\times 2$ 
    matrices), so the task degenerates.
\end{itemize}

For the sub-algebra experiment (Exp.~7), the task is restricted to 
$\fksu(3) \subset \fksu(4)$: inputs lie in the 8-dimensional subspace of 
the 15-dimensional $\fksu(4)$ algebra.

% --- 2.4 ---
\subsection{Architecture}
\label{sec:architecture}

The models compared throughout the campaign are:

\begin{description}[nosep]
  \item[\AlgNet{}:] Algebraic layer (Eq.~\ref{eq:algebraic_layer}) 
    $\to$ classifier head. The algebraic layer has $\ngen$ outputs. 
    The head varies across experiments (linear, 1-layer ReLU, 
    2-layer ReLU).
  \item[MLP:] Standard fully-connected network with matched 
    architecture (same hidden dimension and depth as the \AlgNet{} 
    head, with unconstrained first layer).
  \item[\GMLFixed{}:] Oracle baseline with $\theta = I$ frozen---the 
    algebraic layer is an exact coordinate extractor. Only the 
    head is trained.
\end{description}

All models are trained with Adam~\cite{kingma2015adam} 
($\beta_1 = 0.9$, $\beta_2 = 0.999$, lr $= 10^{-3}$) unless 
otherwise stated, with gradient clipping at $\|\nabla\|_{\max} = 1.0$.


% ========================================================================
% 3. EXPERIMENTS AND RESULTS
% ========================================================================
\section{Experiments and Results}
\label{sec:experiments}

We organize the 14 experiments into three phases that reflect the logical 
progression of the campaign.


% --- 3.1 ---
\subsection{Phase I: Manifold matters (Exp.~1--3)}
\label{sec:phase1}

\paragraph{Exp.~1: Killing-regularized MLP.}
The simplest test: interpret the rows of a standard linear layer 
$W \in \RR^{8 \times 9}$ as $3\times 3$ matrices and add $\TK$ as 
regularization ($\lambda_K = 0.05$, warmup over 50 epochs). 
\textbf{Result}: $\TK$ remains stuck at $\sim 26$ throughout 400 epochs 
of training, and accuracy matches the unregularized MLP ($\sim 65\%$). 
The global minimum $\TK = 0$ is unreachable because the weight rows are 
general real matrices, not hermitian traceless elements of $\fksu(3)$.

\paragraph{Exp.~2: Killing pretraining.}
Two-phase approach: (1)~minimize $\TK$ only for 200 epochs, 
(2)~fine-tune on the task. \textbf{Result}: $\TK$ partially decreases 
during pretraining but never approaches 0. Fine-tuning shows no 
significant improvement over random initialization. This definitively 
closes the unconstrained approach.

\paragraph{Exp.~3: Algebraic layer.}
The fix: use the algebraic parameterization 
$G_a = \sum_k \theta_{ak} T_k$ (Eq.~\ref{eq:algebraic_layer}). 
With $\lambda_K = 2.0$:

\begin{table}[H]
\centering
\caption{Exp.~3 results. The algebraic parameterization enables 
Killing crystallization ($\TK \to 0$).}
\label{tab:exp3}
\begin{tabular}{lccc}
\toprule
Model & Val Acc & $\TK$ final & Crystallized \\
\midrule
MLP        & $\sim$66\% & ---    & --- \\
\AlgNet{}-Free   & $\sim$66\% & drifts & \texttimes \\
\AlgNet{}-TK     & $\sim$68\% & \textbf{0.0000} & \checkmark \\
\AlgNet{}-Full   & $\sim$68\% & \textbf{0.0000} & \checkmark \\
\bottomrule
\end{tabular}
\end{table}

\begin{observation}[G1]
The algebraic parameterization $G_a = \sum_k \theta_{ak} T_k$ is a 
\textbf{necessary condition} for $\TK \to 0$. Without it, the Killing 
minimum is unreachable regardless of optimization strategy.
\end{observation}


% --- 3.2 ---
\subsection{Phase II: Crystallization and task integration (Exp.~4--5)}
\label{sec:phase2}

\paragraph{Exp.~4: $SU(3)$-symmetric cubic task.}
Inputs are now genuine Gell-Mann coordinates $x \in \RR^8$ of 
$\fksu(3)$ elements. An oracle (\GMLFixed{}, $\theta = I$) provides an 
upper bound.

\begin{table}[H]
\centering
\caption{Exp.~4 results on the genuine $SU(3)$-symmetric task.}
\label{tab:exp4}
\begin{tabular}{lccc}
\toprule
Model & Best Val Acc & $\TK$ final & Active Generators \\
\midrule
MLP            & 75.7\% & ---    & --- \\
\AlgNet{}-Free & 75.9\% & drifts & rank-degenerate \\
\AlgNet{}-TK   & \textbf{81.6\%} & 0.0002 & 5 \\
\GMLFixed{}     & 82.2\% & 0.0000 & --- \\
\bottomrule
\end{tabular}
\end{table}

\AlgNet{}-TK reaches 81.6\%, only 0.55\% below the oracle and +5.9\% 
above the unconstrained MLP. Spontaneous symmetry refinement is observed: 
both \AlgNet{}-TK and the oracle converge to 5 active generators, 
corresponding to an $\fksu(2) \subset \fksu(3)$ subalgebra.

\begin{observation}[G2]
$\lambda_K \TK$ reliably drives crystallization when the manifold is 
correct, yielding $\theta \in O(\ngen)$---a genuine Lie algebra basis.
\end{observation}

\begin{observation}[G5]
Emergent representations differ from canonical bases: high 
$\|\theta\theta^T - I\|_F$ coexists with $\TK \approx 0$. The network 
discovers a rotated representation; the Killing metric is the invariant, 
not the matrix form.
\end{observation}

\paragraph{Exp.~5: Scaling study (Parts A--E).}
Five sub-experiments address: (A)~gated rank discovery, (B)~$d$-sweep 
($d = 2$--$8$), (C)~$\lambda_{\mathrm{gate}}$ sweep, (D)~extended 
epochs, and (E)~normalized $\lambda_K$.

Key results from Part B ($d$-sweep):

\begin{table}[H]
\centering
\caption{Exp.~5B: algebraic advantage $\Delta = \text{Acc}(\AlgNet) 
- \text{Acc}(\text{MLP})$ as a function of algebra dimension $d$.}
\label{tab:exp5b}
\begin{tabular}{ccccc}
\toprule
$d$ & $\ngen$ & \AlgNet{}-TK & MLP & $\Delta$ \\
\midrule
3 & 8  & 0.719 & 0.725 & $-0.007$ \\
4 & 15 & 0.629 & 0.623 & $+0.006$ \\
5 & 24 & 0.575 & 0.585 & $-0.010$ \\
6 & 35 & 0.544 & 0.525 & $+0.019$ \\
8 & 63 & 0.525 & 0.504 & $+0.021$ \\
\bottomrule
\end{tabular}
\end{table}

Part E introduces the normalized $\lambda_K^{\mathrm{eff}}$ 
(Eq.~\ref{eq:lambda_norm}), resolving a competition effect at $d = 4$ 
where unnormalized $\lambda_K = 2.0$ overwhelms the task gradient.

\begin{observation}[G3]
The effective Killing gradient scales as $O(\lambda_K \cdot \ngen)$; 
per-generator normalization 
$\lambda_K^{\mathrm{eff}} = 2.0 \cdot 8/(d^2-1)$ is required for 
balanced behavior across dimensions.
\end{observation}

\begin{observation}[G4]
The group-Lasso gate fails to prune generators when the task uses all 
generators symmetrically (the $\fksu(d)$ cubic invariant). It acts as 
beneficial uniform shrinkage instead.
\end{observation}


% --- 3.3 ---
\subsection{Phase III: Manifold confinement (Exp.~6)}
\label{sec:phase3}

Without $\Tortho$, the matrix $\theta$ drifts off the orthogonal 
manifold during training. At $d = 5$, $\Tortho$ reaches 53 (Exp.~5E), 
correlating with crystallization failure. Exp.~6 introduces a 3-phase 
annealing protocol inspired by lattice Monte~Carlo verification of 
exceptional algebras~\cite{moreno2026experimental}:

\begin{table}[H]
\centering
\caption{3-phase annealing protocol for Exp.~6.}
\label{tab:annealing}
\begin{tabular}{llcccc}
\toprule
Phase & Duration & $\lambda_K$ & $\lambda_{\mathrm{ortho}}$ & Head & Analog \\
\midrule
1 --- Task learning & 600 ep & $\lambda_K^{\mathrm{eff}}$ & 0.1 & free & $\gamma_1 = 50$ \\
2 --- Crystallization & 200 ep & $5\lambda_K^{\mathrm{eff}}$ & 2.0 & frozen & $\gamma_2 = 200$ \\
3 --- Fine-tuning & 200 ep & $\lambda_K^{\mathrm{eff}}$ & 0.5 & free & release \\
\bottomrule
\end{tabular}
\end{table}

Results for $d = 4$ ($\fksu(4)$, $\ngen = 15$):

\begin{table}[H]
\centering
\caption{Exp.~6 results at $d = 4$. OrthoAnn beats the oracle.}
\label{tab:exp6_d4}
\begin{tabular}{lcccc}
\toprule
Model & Best Acc & $\Delta$(MLP) & $\TK$ final & $\Tortho$ \\
\midrule
\GMLFixed{} & 0.624 & $-0.001$ & --- & --- \\
MLP & 0.625 & 0 & --- & --- \\
\AlgNet{}-NoOrtho & 0.631 & $+0.007$ & 0.0006 & \textbf{12.11} \\
\AlgNet{}-Ortho1 & 0.631 & $+0.006$ & 0.0005 & 0.02 \\
\AlgNet{}-OrthoAnn & \textbf{0.646} & $+\mathbf{0.021}$ & 0.0005 & 0.02 \\
\bottomrule
\end{tabular}
\end{table}

At $d = 5$, the annealing protocol \emph{hurts} ($-0.021$ vs MLP) 
because Phase~1 reaches only $\sim 56.5\%$ accuracy (barely above 
chance) before the head is frozen, locking in an underdeveloped 
solution.

\begin{observation}[G7]
$\Tortho = \|\theta\theta^T - I\|_F^2$ is necessary to prevent 
catastrophic manifold drift at $d \geq 4$. Without it, 
$\Tortho: 12 \to 53$ from $d = 4 \to 5$.
\end{observation}

\begin{observation}[G8]
3-phase annealing is beneficial only when Phase~1 achieves 
val~acc $>$ (chance $+ 5\%$). For immature representations, 
head-freezing locks in an underdeveloped solution.
\end{observation}

\begin{observation}[G9]
Killing well depth $\lambda_K$ outperforms phase scheduling at high 
$d$. For $d \geq 5$, single-phase training with 
$\lambda_K \in [1.0, 1.5] \times \lambda_K^{\mathrm{eff}}$ is 
strictly better than 3-phase annealing.
\end{observation}


% --- 3.4 ---
\subsection{Phase IV: Sub-algebra gate (Exp.~7)}
\label{sec:phase4}

The group-Lasso failed on the full $\fksu(3)$ cubic task (G4) because 
all generators contribute equally. Exp.~7 tests it on a task with 
genuine sub-algebra structure: the $\fksu(3)$ cubic invariant evaluated 
on $\fksu(3) \subset \fksu(4)$.

\begin{table}[H]
\centering
\caption{Exp.~7: Sub-algebra recovery $\fksu(3) \subset \fksu(4)$.
AlignScore measures alignment with the true 8D signal subspace 
(1.000 = perfect).}
\label{tab:exp7}
\begin{tabular}{lcccc}
\toprule
Model & Best Acc & AlignScore & $s$-pattern \\
\midrule
MLP & 0.901 & --- & --- \\
\AlgNet{}-NoGate & 0.699 & $0.875 \to 0.750$ & unstable \\
\AlgNet{}-Gate-lo ($\lambda_g = 0.05$) & 0.763 & \textbf{1.000} & stable \\
\AlgNet{}-Gate-med ($\lambda_g = 0.10$) & 0.791 & \textbf{1.000} & stable \\
Oracle & 0.855 & 1.000 & --- \\
\bottomrule
\end{tabular}
\end{table}

\begin{observation}[G10]
With a true sub-algebra task, the group-Lasso gate achieves 
\textbf{perfect sub-algebra recovery} (AlignScore $= 1.000$) from 
epoch 300, stably for 1200+ epochs.
\end{observation}

\begin{observation}[G11]
Stronger $\lambda_{\mathrm{gate}}$ monotonically improves accuracy: 
the sparsity pressure is informative (noise suppression), not harmful 
(over-regularization).
\end{observation}


% --- 3.5 ---
\subsection{Phase V: Data efficiency and non-compact algebras (Exp.~8--9)}
\label{sec:phase5}

\paragraph{Exp.~8: Data efficiency sweep.}
Training set size $n_{\mathrm{train}} \in \{100, 300, 1000, 3000, 6000\}$, 
$\fksu(4)$ cubic task, 3 seeds per cell.

\begin{table}[H]
\centering
\caption{Exp.~8: Data efficiency. \AlgNet{} wins at 4/5 data sizes.}
\label{tab:exp8}
\begin{tabular}{ccccc}
\toprule
$n_{\mathrm{train}}$ & MLP (mean$\pm$std) & \AlgNet{} (mean$\pm$std) & $\Delta$ & Winner \\
\midrule
100   & $0.514 \pm 0.006$ & $0.531 \pm 0.007$ & $+0.017$ & \AlgNet{} \\
300   & $0.530 \pm 0.005$ & $0.543 \pm 0.006$ & $+0.014$ & \AlgNet{} \\
1000  & $0.574 \pm 0.006$ & $0.571 \pm 0.008$ & $-0.003$ & MLP (n.s.) \\
3000  & $0.607 \pm 0.014$ & $0.624 \pm \mathbf{0.003}$ & $+0.018$ & \AlgNet{} \\
6000  & $0.632 \pm 0.007$ & $0.637 \pm 0.007$ & $+0.004$ & \AlgNet{} \\
\bottomrule
\end{tabular}
\end{table}

The algebraic advantage is not monotonically increasing with scarcity 
(peak $\Delta$ at $n = 3000$). The key finding is the 
$4.6\times$ lower variance of \AlgNet{} at $n = 3000$: the algebraic 
prior acts primarily as a \emph{stability regularizer}.

\begin{observation}[G12]
\AlgNet{} wins at 4/5 data sizes, but the advantage is not 
monotonically increasing with data scarcity. Peak $\Delta$ at 
$n = 3000$; \AlgNet{} shows $4.6\times$ lower variance---the 
algebraic prior is a stability regularizer.
\end{observation}

\paragraph{Exp.~9: Non-compact algebra $\fksl(3,\RR)$.}
The $\fksl(3,\RR)$ algebra has an \textbf{indefinite} Killing form 
(5 positive, 3 negative eigenvalues). Three variants are tested: 
MLP (no prior), \AlgNet{}-NC (correct non-compact Killing target), and 
\AlgNet{}-Compact (wrong compact target $K \propto -I$).

\begin{table}[H]
\centering
\caption{Exp.~9: Non-compact $\fksl(3,\RR)$ task (3 seeds).}
\label{tab:exp9}
\begin{tabular}{lccc}
\toprule
Model & Mean Acc & std & vs MLP \\
\midrule
MLP & 0.753 & 0.011 & --- \\
\AlgNet{}-NC (correct) & \textbf{0.831} & \textbf{0.005} & $+0.079$ ($p = 0.004$) \\
Oracle & 0.822 & --- & $+0.069$ \\
\AlgNet{}-Compact (wrong) & 0.671 & \textbf{0.112} & $-0.082$ \\
\bottomrule
\end{tabular}
\end{table}

The correct non-compact prior not only improves over MLP 
($+7.9\%$, $p = 0.004$) but \textbf{exceeds the oracle} (0.831 vs 
0.822). The wrong compact prior is catastrophic: bimodal collapse with 
$10\times$ the MLP variance.

\begin{observation}[G13]
Prior consistency determines everything for non-compact algebras. 
Correct NC prior: $+7.9\%$ vs MLP ($p = 0.004$), exceeds oracle. 
Wrong compact prior: bimodal collapse (std $10\times$ MLP). $\Tortho$ 
alone recovers the correct K-signature.
\end{observation}


% --- 3.6 ---
\subsection{Phase VI: The decoder bottleneck (Exp.~10, 12, 14)}
\label{sec:phase6}

The pivotal discovery of the campaign: the algebraic layer 
\emph{always crystallizes correctly}, but whether that helps the task 
depends on the decoder architecture.

\paragraph{Exp.~10: Stacked algebraic layers ($d = 3$, depth 1--3).}
With a thin linear head $\text{Linear}(\ngen, 2)$:

\begin{table}[H]
\centering
\caption{Exp.~10: Depth sweep with thin (linear) head.}
\label{tab:exp10}
\begin{tabular}{lcccc}
\toprule
Depth & \AlgNet{} & MLP & $\Delta$ & $p$-value \\
\midrule
1 & 0.707 & 0.744 & $-0.037$ & 0.001*** \\
2 & 0.723 & 0.764 & $-0.041$ & 0.006** \\
3 & 0.742 & 0.771 & $-0.029$ & 0.26 (n.s.) \\
\bottomrule
\end{tabular}
\end{table}

\AlgNet{} trails MLP at all depths, but the gap narrows from 
$-3.7\%$ (depth 1) to $-2.9\%$ (depth 3, non-significant). Crucially, 
\textbf{the Killing signature is exact at all depths and seeds}: the 
algebraic constraint is working, but the thin head cannot decode the 
information.

\begin{observation}[G14]
With a thin head Linear($\ngen$, 2), the algebraic constraint is a 
bottleneck: \AlgNet{} trails MLP at depths 1,2 (significant). Depth 
helps \AlgNet{} more ($+11.9\%$) than MLP ($+7.2\%$).
\end{observation}

\paragraph{Exp.~12: Fine $d$-sweep ($d = 2$--$12$, 10 seeds).}
With the thin head, no positive $\Delta(d)$ peak exists for any $d$. 
Significant \AlgNet{} deficit at $d = 3, 4$; both models hit the 
accuracy floor at $d \geq 5$.

\begin{observation}[G16]
Fine $d$-sweep with thin head: no positive $\Delta(d)$ peak exists in 
$d = 2..12$. Significant \AlgNet{} deficit only at $d = 3, 4$; all 
models hit the accuracy floor at $d \geq 5$.
\end{observation}

\paragraph{Exp.~14: Head width sweep (1L and 2L heads, $h \in \{2, 4, 8, 16, 32\}$, 10 seeds).}
This is the decisive experiment:

\begin{table}[H]
\centering
\caption{Exp.~14: Head architecture sweep. The 2-layer head with 
$h = 32 \approx 4 \times \ngen$ unlocks the algebraic advantage.}
\label{tab:exp14}
\begin{tabular}{llcccc}
\toprule
Head & $h$ & \AlgNet{} & MLP & $\Delta$ & Seeds $+$ \\
\midrule
1L & 2  & 0.698 & 0.737 & $-0.039$ & 0/10 \\
1L & 8  & 0.714 & 0.748 & $-0.034$ & 0/10 \\
1L & 32 & 0.719 & 0.746 & $-0.027$ & 2/10 \\
\midrule
2L & 2  & 0.699 & 0.741 & $-0.042$ & 0/10 \\
2L & 8  & 0.740 & 0.753 & $-0.013$ & 2/10 \\
2L & 16 & 0.758 & 0.747 & $+0.011$ & 7/10 \\
2L & 32 & \textbf{0.791} & 0.725 & $+\mathbf{0.067}$ & \textbf{10/10} \\
\bottomrule
\end{tabular}
\end{table}

\begin{observation}[G18 --- the key result]
\textbf{The Stiefel bottleneck is a decoder problem.} 1L head: \AlgNet{} 
never wins. 2L head: crossover at $h^* = 32 \approx 4 \times \ngen$. 
\AlgNet{}-2L-h32 achieves 0.791 (record, $+6.65\%$ vs MLP, 10/10 seeds 
positive). G18 retroactively explains all prior negative results (G14, G16).
\end{observation}

The resolution is clear: the $O(\ngen)$-constrained features from the 
algebraic layer are \textbf{linearly indecodable} for the cubic invariant 
$\Tr(M^3)$, which is a degree-3 polynomial. A single linear layer 
cannot express this. Once the head has depth $\geq 2$ and width 
$h \geq 4\ngen$, the algebraic prior provides a decisive advantage 
because it constrains the representation to the exact symmetry manifold, 
while the nonlinear head can express the necessary polynomial function.


% --- 3.7 ---
\subsection{Phase VII: Optimizer dynamics (Exp.~11, 13)}
\label{sec:phase7}

\paragraph{Exp.~11: Fractal dimension.}
Using the Grassberger--Procaccia algorithm~\cite{grassberger1983} to 
compute the correlation dimension $D_2$ of optimizer trajectories in 
$\theta$-space ($\fksu(3)$, 5 seeds, 3000 epochs):

\begin{table}[H]
\centering
\caption{Exp.~11: Correlation dimension $D_2$ of optimizer 
trajectories in different training phases.}
\label{tab:exp11}
\begin{tabular}{lccc}
\toprule
Phase & \AlgNet{} $D_2$ & MLP $D_2$ & Interpretation \\
\midrule
Full trajectory   & $< 1$ & $< 1$ & Both fractal \\
Convergence phase & \textbf{0.99} & 1.41  & \AlgNet{} more confined \\
\bottomrule
\end{tabular}
\end{table}

\begin{observation}[G15]
Optimizer trajectories are fractal objects ($D_2 < 1$). In the 
convergence phase, $\Tortho$ confines \AlgNet{} to a 
lower-dimensional attractor ($D_2 = 0.99$) vs MLP ($D_2 = 1.41$), 
confirming Stiefel geometry dynamically.
\end{observation}

\paragraph{Exp.~13: Lyapunov exponents.}
Finite-time Lyapunov exponents~\cite{wolf1985determining} of the Adam 
dynamics in the full parameter space $(\theta, m, v)$ 
($\fksu(3)$, 5 seeds, 3000 epochs):

\begin{table}[H]
\centering
\caption{Exp.~13: Lyapunov exponents of Adam dynamics.}
\label{tab:exp13}
\begin{tabular}{lcc}
\toprule
Exponent & \AlgNet{} & MLP \\
\midrule
$\lambda_\theta$ (parameter space)   & $\approx +0.20$/step & $\approx +0.20$/step \\
$\lambda_{\mathrm{ext}}$ (extended phase space) & \textbf{0.215} & 0.232 \\
\bottomrule
\end{tabular}
\end{table}

\begin{observation}[G17]
Both models exhibit deterministic chaos ($\lambda_\theta \approx +0.20$/step 
throughout 3000 epochs). \AlgNet{}'s extended-phase-space exponent 
($\lambda_{\mathrm{ext}}$) is systematically lower: $\Tortho$ regularizes 
Adam momentum buffers. No late-phase collapse---training lives on a 
persistent strange attractor.
\end{observation}


% ========================================================================
% 4. GENERALIZATIONS
% ========================================================================
\section{Consolidated Generalizations}
\label{sec:generalizations}

The 18 empirical generalizations are organized into seven thematic groups.

\subsection{I. Manifold and Parameterization}

\noindent\textbf{G1.} The algebraic parameterization $G_a = \sum_k \theta_{ak} T_k$ is a 
\emph{necessary condition} for $\TK \to 0$. Without it, the Killing minimum is 
unreachable. [\emph{Exp.~1--3}]

\noindent\textbf{G5.} Emergent representations differ from canonical bases: high 
$\|\theta\theta^T - I\|_F$ coexists with $\TK \approx 0$. The network finds a 
rotated representation---the Killing metric is the invariant, not the matrix form. 
[\emph{Exp.~4--5}]

\subsection{II. Killing Dynamics and Crystallization}

\noindent\textbf{G2.} $\lambda_K \TK$ reliably drives crystallization ($\TK \to 0$) when the 
manifold is correct, yielding $\theta \in O(\ngen)$---a genuine Lie algebra basis. 
[\emph{Exp.~3--4}]

\noindent\textbf{G3.} The effective Killing gradient scales as $O(\lambda_K \cdot \ngen)$; 
per-generator normalization $\lambda_K^{\mathrm{eff}} = 2.0 \cdot 8/(d^2-1)$ is 
required for balanced behavior across dimensions. [\emph{Exp.~5}]

\noindent\textbf{G9.} Killing well depth $\lambda_K$ outperforms phase scheduling at high $d$. 
For $d \geq 5$, single-phase training with 
$\lambda_K \in [1.0, 1.5] \times \lambda_K^{\mathrm{eff}}$ is strictly better than 
3-phase annealing. [\emph{Exp.~6}]

\subsection{III. Manifold Confinement and Annealing}

\noindent\textbf{G7.} $\Tortho = \|\theta\theta^T - I\|_F^2$ is the neural analog of 
$\gamma \cdot \Tclosure$ in lattice verification. Without it, $\theta$ drifts off 
$O(\ngen)$ catastrophically ($\Tortho$: $12 \to 53$ from $d = 4 \to 5$). 
[\emph{Exp.~5--6}]

\noindent\textbf{G8.} 3-phase annealing is beneficial only when Phase~1 achieves 
val~acc $>$ (chance $+ 5\%$). For immature representations, head-freezing locks in 
an underdeveloped solution. [\emph{Exp.~6}]

\subsection{IV. Sub-Algebra Discovery and Sparsity}

\noindent\textbf{G4.} Group-Lasso gate fails to prune generators when the task uses all generators 
symmetrically. It acts as beneficial uniform shrinkage instead. [\emph{Exp.~5A,C}]

\noindent\textbf{G10.} With a true sub-algebra task ($\fksu(3) \subset \fksu(4)$), the 
group-Lasso gate achieves \emph{perfect} sub-algebra recovery (AlignScore $= 1.000$) 
from epoch 300, stably for 1200+ epochs. [\emph{Exp.~7}]

\noindent\textbf{G11.} Stronger $\lambda_{\mathrm{gate}}$ monotonically improves accuracy: the 
sparsity pressure is informative (noise suppression), not harmful (over-regularization). 
[\emph{Exp.~7}]

\subsection{V. Data Efficiency and Non-Compact Algebras}

\noindent\textbf{G12.} \AlgNet{} wins at 4/5 data sizes, but the advantage is not monotonically 
increasing with data scarcity. Peak $\Delta$ at $n = 3000$; \AlgNet{} shows $4.6\times$ 
lower variance---the algebraic prior is a stability regularizer. [\emph{Exp.~8}]

\noindent\textbf{G13.} Prior consistency determines everything for non-compact algebras. Correct 
NC prior: $+7.9\%$ vs MLP ($p = 0.004$), exceeds oracle. Wrong compact prior: bimodal 
collapse (std $10\times$ MLP). [\emph{Exp.~9}]

\subsection{VI. Architecture: The Decoder Bottleneck}

\noindent\textbf{G14.} With a thin head Linear($\ngen$, 2), the algebraic constraint is a 
bottleneck: \AlgNet{} trails MLP at depths 1, 2 (significant), closing to n.s.\ at 
depth 3. Depth helps \AlgNet{} more ($+11.9\%$) than MLP ($+7.2\%$). [\emph{Exp.~10}]

\noindent\textbf{G16.} Fine $d$-sweep with thin head: no positive $\Delta(d)$ peak exists in 
$d = 2..12$. Significant \AlgNet{} deficit only at $d = 3, 4$; all models hit accuracy 
floor at $d \geq 5$. [\emph{Exp.~12}]

\noindent\textbf{G18.} \textbf{The Stiefel bottleneck is a decoder problem.} 1L head: \AlgNet{} 
NEVER wins. 2L head: crossover at $h^* = 32 \approx 4 \times \ngen$. 
\AlgNet{}-2L-h32 achieves 0.791 (record, $+6.65\%$ vs MLP, 10/10 seeds positive). 
G18 retroactively explains all prior negative results. [\emph{Exp.~14}]

\subsection{VII. Optimizer Dynamics: Chaos and Fractals}

\noindent\textbf{G15.} Optimizer trajectories are fractal objects ($D_2 < 1$). In the 
convergence phase, $\Tortho$ confines \AlgNet{} to a lower-dimensional attractor 
($D_2 = 0.99$) vs MLP ($D_2 = 1.41$). [\emph{Exp.~11}]

\noindent\textbf{G17.} Both models exhibit deterministic chaos ($\lambda_\theta \approx +0.20$/step 
throughout 3000 epochs). \AlgNet{}'s extended-phase-space exponent 
($\lambda_{\mathrm{ext}}$) is systematically lower: $\Tortho$ regularizes Adam 
momentum buffers. [\emph{Exp.~13}]


% ========================================================================
% 5. DISCUSSION
% ========================================================================
\section{Discussion}
\label{sec:discussion}

\subsection{The decoder bottleneck: a unifying narrative}
\label{sec:bottleneck}

The single most important finding of this campaign is that the 
\emph{decoder architecture}, not the algebraic layer, determines whether 
the Killing-driven inductive bias translates into a task advantage. This 
resolves an apparent paradox: the algebraic layer always crystallizes 
correctly (Killing signature exact at all depths, dimensions, and seeds), 
yet early experiments (Exp.~10, 12) showed \AlgNet{} trailing MLP.

The resolution is architectural. The algebraic layer constrains its 
outputs to $O(\ngen)$-rotated coordinates---a linear isometry of the 
original Gell-Mann basis. The cubic invariant $\Tr(M^3)$ is a 
\textbf{degree-3 polynomial} of these coordinates. A single linear 
projection $\text{Linear}(\ngen, 2)$ cannot express this function, 
creating an information bottleneck where the algebraic constraint 
\emph{reduces} the effective hypothesis class below what the MLP 
(unconstrained first layer, same head) can achieve.

Once the head provides sufficient nonlinear capacity---specifically, 
depth $\geq 2$ and width $h \geq 4\ngen$---the algebraic constraint 
becomes purely beneficial: it eliminates the $(d^2-1)^2 - \ngen(\ngen-1)/2$ 
off-manifold degrees of freedom that the MLP must learn to ignore, 
acting as an exact symmetry regularizer. The crossover is sharp: 
at $h = 16$ (2L), 7/10 seeds are positive; at $h = 32$, 10/10 seeds 
are positive with $\Delta = +6.65\%$.

This finding has a direct analog in the physical verification: the 
closure tension $\gamma \cdot \Tclosure$ in the lattice Monte~Carlo is 
the confinement mechanism, but it only produces useful results when the 
loss landscape has a basin structure compatible with the target 
algebra~\cite{moreno2026experimental}. In the neural setting, the 
``basin structure'' is determined by the decoder: a sufficiently rich 
head creates a loss landscape where the algebraic constraint's basin 
aligns with the task objective.


\subsection{Cross-experiment insights}
\label{sec:cross_insights}

\paragraph{The algebraic layer always crystallizes correctly.}
Across 14 experiments---every dimension ($d = 2..12$), every depth 
(1--3 layers), compact and non-compact algebras, 3-seed to 10-seed 
averages---the Killing signature converges to the reference value 
whenever $\TK$ is minimized. The mechanism is robust. The question 
was never ``does it crystallize?'' but ``does crystallization help 
the task?''

\paragraph{Prior consistency matters more than prior existence.}
Exp.~9 (G13) demonstrated that the wrong algebraic prior (compact 
$\TK$ on a non-compact algebra) is worse than no prior at all---catastrophically 
so (bimodal collapse, $10\times$ variance). The correct non-compact 
prior exceeds even the oracle. In machine learning terms: the inductive 
bias must match the data symmetry; a mismatched structural prior creates 
a dangerously multimodal loss landscape.

\paragraph{The SS-PEG physical analogy holds precisely.}
The neural $\Tortho$ is the functional equivalent of the lattice 
$\gamma \cdot \Tclosure$. Both systems require manifold confinement to 
prevent drift, both benefit from annealing protocols with phase 
transitions, and both produce emergent algebraic structure from 
variational pressure alone. The optimizer dynamics even share 
qualitative features: fractal attractors, positive Lyapunov exponents, 
and phase-dependent dimensionality.


\subsection{Design rules for practitioners}
\label{sec:design_rules}

Based on the 18 generalizations, we distill these rules for deploying 
algebraic inductive biases in neural networks:

\begin{table}[H]
\centering
\caption{Design rules derived from G1--G18.}
\label{tab:design_rules}
\begin{tabularx}{\textwidth}{lXc}
\toprule
Rule & Description & Source \\
\midrule
Parameterization & Use $G_a = \sum_k \theta_{ak} T_k$, never 
  unconstrained matrices & G1 \\
Confinement & Always include $\Tortho$; omitting it causes catastrophic 
  manifold drift at $d \geq 4$ & G7 \\
$\lambda_K$ scaling & Normalize by $\ngen$: 
  $\lambda_K^{\mathrm{eff}} = \lambda_K^{\mathrm{ref}} 
  \cdot \ngen^{\mathrm{ref}} / \ngen$ & G3 \\
Head architecture & Use a 2-layer$+$ head with width 
  $h \geq 4 \times \ngen$ & G18 \\
Compact algebras & $\lambda_K \in [1.0, 1.5] \times \lambda_K^{\mathrm{eff}}$ 
  in single phase; avoid 3-phase annealing unless Phase~1 
  acc $>$ chance $+ 5\%$ & G8, G9 \\
Non-compact & Use the correct $K_{\mathrm{ref}}$ target; compact $\TK$ 
  is catastrophic & G13 \\
Sub-algebra & Add group-Lasso gate ($\lambda_{\mathrm{gate}} \geq 0.05$) 
  when sub-algebra structure is expected & G10, G11 \\
\bottomrule
\end{tabularx}
\end{table}


\subsection{Connection to the physical framework}
\label{sec:physics}

The experiments in this paper test the machine learning implication of a 
broader theoretical framework~\cite{moreno2026relational} in which the 
Killing metric is the unique geometric object that drives symmetry 
emergence from relational potentials. The key correspondences are:

\begin{itemize}[nosep]
  \item The algebraic layer's parameter $\theta$ corresponds to the 
    relational potential in the theoretical framework.
  \item The Killing tension $\TK$ is the variational functional whose 
    minimization produces algebraic structure.
  \item The confinement term $\Tortho$ is the neural analog of the 
    closure tension $\Tclosure$ used in lattice verification of 
    exceptional algebras.
  \item The group-Lasso gate corresponds to a sparsity-driven selection 
    mechanism that identifies the relevant subalgebra.
  \item The decoder bottleneck has no direct physical analog but 
    determines the practical utility of the algebraic prior in the 
    ML setting.
\end{itemize}

The experimental campaign establishes that the Killing metric---a purely 
geometric object defined by the adjoint representation---contains 
\emph{all} the information needed to reconstruct a Lie algebra basis 
within a neural network layer. The algebra is determined by its own 
geometry: a self-referential structure consistent with the 
potential $\to$ cycle $\to$ invariant chain described in the theoretical 
framework.


% ========================================================================
% 6. RECORD ACCURACY SUMMARY
% ========================================================================
\section{Record Accuracy Summary}
\label{sec:records}

\begin{table}[H]
\centering
\caption{Best validation accuracies achieved across the campaign.}
\label{tab:records}
\begin{tabular}{lcc}
\toprule
Configuration & Best Val Acc & Experiment \\
\midrule
\AlgNet{}-2L-h32, $\fksu(3)$ cubic    & \textbf{0.791} & Exp.~14 \\
\AlgNet{}-NC, $\fksl(3,\RR)$ cubic    & \textbf{0.831} & Exp.~9 \\
\AlgNet{}-TK, $\fksu(3)$ GML-input    & 0.816          & Exp.~4 \\
\AlgNet{}-Gate-med, $\fksu(3) \subset \fksu(4)$ & 0.791 & Exp.~7 \\
\AlgNet{}-OrthoAnn, $\fksu(4)$ cubic  & 0.646          & Exp.~6 \\
\bottomrule
\end{tabular}
\end{table}


% ========================================================================
% 7. CONCLUSION
% ========================================================================
\section{Conclusion}
\label{sec:conclusion}

This 14-experiment campaign demonstrates that \textbf{Lie algebra 
structure can emerge as an inductive bias in neural networks through 
variational pressure alone}. The mechanism is reliable: crystallization 
succeeds across all tested algebras and configurations whenever the 
algebraic parameterization is used. The critical insight---that the 
algebraic layer is an encoder whose value is gated by the 
decoder---resolves all apparent contradictions in the experimental 
record and provides clear design principles for future applications.

The main findings are:
\begin{enumerate}[nosep]
  \item The algebraic parameterization $G_a = \sum_k \theta_{ak} T_k$ is 
    necessary and sufficient for Killing crystallization (G1, G2).
  \item Manifold confinement via $\Tortho$ is required at $d \geq 4$ (G7).
  \item The Killing coefficient must be normalized by $\ngen$ for 
    cross-dimensional stability (G3).
  \item A group-Lasso gate achieves perfect sub-algebra recovery when 
    sub-algebra structure exists in the data (G10).
  \item The algebraic prior is most valuable for non-compact algebras when 
    the correct Killing target is used (G13).
  \item \textbf{The decoder head is the decisive variable}: a 2-layer head 
    with $h \geq 4\ngen$ unlocks the full algebraic advantage, yielding 
    $+6.65\%$ over MLP with 10/10 seeds positive (G18).
  \item Training dynamics on the algebraic manifold exhibit deterministic 
    chaos with lower fractal dimension (G15, G17), indicating that 
    $\Tortho$ provides geometric regularization of the optimizer 
    trajectory.
\end{enumerate}

These results establish concrete, actionable design rules for deploying 
Lie algebra inductive biases in neural networks and validate the 
Killing metric as the sole variational driver of symmetry emergence in 
learned representations.


% ========================================================================
% ACKNOWLEDGMENTS
% ========================================================================
\section*{Acknowledgments}

Computations were performed on a single NVIDIA RTX 5070~Ti 
(17.1~GB VRAM) with CUDA~13 and Python~3.14 (free-threaded build).


% ========================================================================
% APPENDIX
% ========================================================================
\appendix

\section{Notation}
\label{app:notation}

\begin{table}[H]
\centering
\begin{tabular}{ll}
\toprule
Symbol & Meaning \\
\midrule
$\fkg$ & Lie algebra \\
$G_a$, $T_a$ & Generators of the algebra \\
$\ngen$ & Number of generators \\
$d$ & Dimension of the matrix representation \\
$K_{ab}$ & Killing metric: $\Tr(\mathrm{ad}_a \circ \mathrm{ad}_b)$ \\
$\kappa$ & Generator norm: $\Tr(T_a^2)/d$ \\
$f_{abc}$ & Structure constants: $[T_a, T_b] = i f_{abc} T_c$ \\
$\hdual$ & Dual Coxeter number \\
$I_R$ & Dynkin index of representation $R$ \\
$C_\lambda$ & Casimir eigenvalue \\
$\TK$ & Killing tension \\
$\Tortho$ & Orthogonality tension: $\|\theta\theta^T - I\|_F^2$ \\
$\Tclosure$ & Closure tension (lattice analog) \\
$\Tgate$ & Gate penalty: $\sum_a s_a$ \\
$\lambda_K^{\mathrm{eff}}$ & Normalized Killing coefficient \\
$D_2$ & Grassberger--Procaccia correlation dimension \\
$\lambda_\theta$ & Finite-time Lyapunov exponent \\
\bottomrule
\end{tabular}
\end{table}


\section{Complete experiment index}
\label{app:experiments}

\begin{table}[H]
\centering
\small
\begin{tabular}{clll}
\toprule
Exp. & Description & Key Finding & Generalizations \\
\midrule
1  & Killing-regularized MLP        & Wrong manifold: $\TK$ stuck at 26  & G1 \\
2  & Killing pretraining             & Pretraining cannot fix manifold    & G1 \\
3  & Algebraic layer                 & $\TK \to 0$ confirmed              & G1, G2 \\
4  & $SU(3)$ symmetric task          & $+5.9\%$ vs MLP, oracle gap 0.55\% & G2, G5 \\
5  & Scaling study (A--E)            & $\lambda_K$ normalization, $d$-sweep & G3, G4 \\
6  & $\Tortho$ annealing             & 3-phase protocol, $+2.1\%$ at $d=4$ & G7, G8, G9 \\
7  & Sub-algebra gate                & Perfect recovery, AlignScore=1.0  & G10, G11 \\
8  & Data efficiency sweep           & 4/5 wins, $4.6\times$ lower var   & G12 \\
9  & Non-compact $\fksl(3,\RR)$      & $+7.9\%$, exceeds oracle           & G13 \\
10 & Stacked depth (thin head)       & Bottleneck identified              & G14 \\
11 & Fractal dimension               & $D_2 = 0.99$ vs 1.41              & G15 \\
12 & Fine $d$-sweep (thin head)      & No positive peak at any $d$        & G16 \\
13 & Lyapunov exponents              & $\lambda_\theta \approx +0.20$/step & G17 \\
14 & Head width sweep                & \textbf{$h \geq 4\ngen$ unlocks advantage} & \textbf{G18} \\
\bottomrule
\end{tabular}
\end{table}


% ========================================================================
% REFERENCES
% ========================================================================
\bibliographystyle{unsrt}
\bibliography{references}

\end{document}
